\Needspace{18\baselineskip}
\section{Typed observations and contextual lifting}

\begin{definition}[Observation interface]
Let $J$ be a type of interfaces, and for each $i\in J$ let $F_i$ be a
type of fragments and $C_i$ a type of contexts. An observation interface
consists of a predicate $\mathsf{Allowed}(c)$ for contexts and a relation
\[
 \mathsf{Obs}(c,a,s),\qquad c\in C_i,\ a\in F_i,\ s\in X.
\]
The intended reading is that placing $a$ in $c$ and initializing the
result can produce state $s$.
\end{definition}
\lean{Contextual}{15}{18}{Observations}
The common index $i$ prevents a fragment and a context of incompatible
interfaces from being compared. Observation is a relation rather than a
total function. It can be undefined, and the structure does not require
uniqueness of the produced state. The current code does not yet define it
by graph plugging.

\begin{definition}[Contextual refinement and equivalence]
For $a,b\in F_i$, write $a\preceq_Q^{\mathcal O} b$ if for every
$c\in C_i$ and states $s,t$,
\[
 \mathsf{Allowed}(c)\ \land\ \mathsf{Obs}(c,a,s)\ \land\ \mathsf{Obs}(c,b,t)
 \ \Longrightarrow\ s\preceq_Q t.
\]
Contextual equivalence $a\simeq_Q^{\mathcal O}b$ is the conjunction of
this refinement and the refinement in the opposite direction.
\end{definition}
\lean{Contextual}{26}{28}{Observations.Refines}
\lean{Contextual}{30}{31}{Observations.Equiv}
The quantifiers compare every pair of observations of the two fragments
in the same allowed context. The definition imposes no existence obligation
when either fragment has no observation there. Also, state-refinement
reflexivity alone does not establish contextual reflexivity for this fully
general interface: two observations of the same fragment might give
different states. A concrete, functional initialization relation can supply
the missing property, but it is not assumed here.

\begin{proposition}[Contextual refinement from a lifted relation]
Let $T_i$ be a relation on $F_i$ for each interface $i$, and let $L$ be
a relation on states. Suppose:
\begin{enumerate}
\item Whenever $T_i(a,b)$ and $c$ is allowed, observations
  $\mathsf{Obs}(c,a,s)$ and $\mathsf{Obs}(c,b,t)$ imply $L(s,t)$.
\item $L(s,t)$ implies $s\preceq_Q t$.
\end{enumerate}
Then $T_i(a,b)$ implies $a\preceq_Q^{\mathcal O}b$.
\end{proposition}
\lean{Contextual}{35}{44}{Observations.contextual_of_lifting}
\begin{proof}
Fix $T_i(a,b)$, an allowed $c$, and the two observations. The first
hypothesis, named \texttt{initial\_lift} in Lean, puts their initial states
in $L$. The second hypothesis, named \texttt{sound}, then gives the state
refinement required by the definition of contextual refinement.
\end{proof}
This isolates the logical step of the paper's Proposition~38. Its concrete
proof still has to establish the initial-specimen hypothesis for actual
plugging and initialization. The proposition above neither asserts nor
requires closure under plugging at this abstract level.

\begin{proposition}[Contextual equivalence from a lifted relation]
Under the same initial-lifting hypothesis, if $L(s,t)$ implies
$s\simeq_Q t$, then $T_i(a,b)$ implies $a\simeq_Q^{\mathcal O}b$.
\end{proposition}
\lean{Contextual}{46}{57}{Observations.contextual_equiv_of_lifting}
\begin{proof}
For the forward refinement, apply the preceding proposition using the
forward component of state equivalence. For the reverse refinement, fix
observations of $b$ and $a$ in an allowed context. Put them into
\texttt{initial\_lift} in the order $a,b$, obtaining $L$ from the state
of $a$ to the state of $b$. Use the reverse component of its state
equivalence. This does not require that $T_i$ or $L$ be symmetric.
\end{proof}
